% --- Archon physics formalization source begin ---
% archon:physics
% archon:covers PhyXMiniProblems/problem_phyx_mini_0479.lean
% archon:source-report reports/phyx_mini/problem_phyx_mini_0479.source.json

\chapter{Physics problem phyx\_mini\_0479}
\label{ch:PhyXMiniProblems_problem_phyx_mini_0479}

\paragraph{Problem source.}
\#\# Physical scenario

Analyze the heat engine of figure.It runs at \textbackslash{}(600\textbackslash{},\textbackslash{}mathrm\{rpm\}\textbackslash{}). Assume the gas is monatomic.

\#\# Question

Determine the engine's power output

\#\# Answer choices

- A: A: 320W

- B: B: 340W

- C: C: 360W

- D: D: 380W

\#\# Figure caption (auxiliary; use the image as primary evidence)

The image is a pressure-volume (p-V) diagram showing a thermodynamic cycle. Here's a detailed description:

1. **Axes:**
   - The y-axis represents pressure (p) measured in kilopascals (kPa), ranging from 0 to over 200 kPa.
   - The x-axis represents volume (V) measured in cubic centimeters (cm³), ranging from 0 to 600 cm³.

2. **Diagram Elements:**
   - There are three key points labeled 1, 2, and 3.
   - The cycle is depicted with arrows indicating direction: from point 1 to 2, from 2 to 3, and from 3 back to 1.

3. **Processes:**
   - **1 to 2:** A horizontal line indicating a constant pressure process at around 200 kPa.
   - **2 to 3:** A vertical line indicating a constant volume process at about 600 cm³.
   - **3 to 1:** A curved line reflecting an isothermal process indicated as a 300 K isotherm, leading back to point 1.

4. **Isotherm:**
   - A dashed red line labeled "300 K isotherm" runs roughly parallel to the curve from point 3 to 1, showing a process at constant temperature. 

5. **Text:**
   - The text "300 K isotherm" is positioned near the dashed curve.

This diagram is typical for representing processes in an idealized thermodynamic cycle.

\#\# Dataset metadata

Laws of Thermodynamics; Physical Model Grounding Reasoning, Multi-Formula Reasoning

\paragraph{Recorded answer/context.}
C: C: 360W

\paragraph{Figure/image path.}
/root/proposal\_for\_physic/science-mango/phyx\_mini\_run/phyx\_data/test\_image/479.png

\paragraph{Formalization target.}
create a compiling Lean file with sorry bodies at `PhyXMiniProblems/problem\_phyx\_mini\_0479.lean`. The Lean declarations must preserve the physical quantities, dimensions or dimensional roles, figure labels, governing-law hypotheses, and final relation expressed by this problem.
Use Mathlib/Physlib names found through LeanExplore where available. If a physics API is missing, introduce faithful local abstractions rather than scalar placeholder aliases.

\begin{theorem}[Physics formalization target]
\label{thm:physics:phyx_mini_0479:target}
\lean{PhyXMiniProblems.ProblemPhyXMini0479.problem_phyx_mini_0479}
\uses{def:physics:phyx-mini-0479:phyxminiproblems-problemphyxmini0479-powerinwatts, def:physics:phyx-mini-0479:phyxminiproblems-problemphyxmini0479-idealgasheatenginecycle, def:physics:phyx-mini-0479:phyxminiproblems-problemphyxmini0479-matchesproblemstatement, def:physics:phyx-mini-0479:phyxminiproblems-problemphyxmini0479-matchesprimarypressurevolumefigure, def:physics:phyx-mini-0479:phyxminiproblems-problemphyxmini0479-hasphysicalthermodynamicparameters, def:physics:phyx-mini-0479:phyxminiproblems-problemphyxmini0479-satisfiesmonatomicidealgasenginelaws, def:physics:phyx-mini-0479:phyxminiproblems-problemphyxmini0479-recordedanswerchoice, def:physics:phyx-mini-0479:phyxminiproblems-problemphyxmini0479-roundstodisplayedpower, def:physics:phyx-mini-0479:phyxminiproblems-problemphyxmini0479-isuniqueclosestdisplayedpower}
The assigned autoformalize agent should translate this physics problem into Lean declarations in the covered file, with theorem and lemma proof bodies written as `by sorry`.
\end{theorem}
\begin{proof}
This is an autoformalization task, not a proof task. Produce statements that can later be proved by the physics prover without weakening the physical model.
\end{proof}

% --- Archon declaration topology begin ---
\section{Declaration topology}
\begin{definition}[Volume Quantity]
  \label{def:physics:phyx-mini-0479:phyxminiproblems-problemphyxmini0479-volumequantity}
  \lean{PhyXMiniProblems.ProblemPhyXMini0479.VolumeQuantity}
  A nonnegative physical volume carrying dimension L³
\end{definition}
\begin{proof}
  This is the declared model definition of Volume Quantity.
\end{proof}
\begin{definition}[Rotation Rate Quantity]
  \label{def:physics:phyx-mini-0479:phyxminiproblems-problemphyxmini0479-rotationratequantity}
  \lean{PhyXMiniProblems.ProblemPhyXMini0479.RotationRateQuantity}
  Shaft rotation rate, carrying inverse-time dimension
\end{definition}
\begin{proof}
  This is the declared model definition of Rotation Rate Quantity.
\end{proof}
\begin{definition}[Cycle Rate Quantity]
  \label{def:physics:phyx-mini-0479:phyxminiproblems-problemphyxmini0479-cycleratequantity}
  \lean{PhyXMiniProblems.ProblemPhyXMini0479.CycleRateQuantity}
  Thermodynamic cycle rate, carrying inverse-time dimension
\end{definition}
\begin{proof}
  This is the declared model definition of Cycle Rate Quantity.
\end{proof}
\begin{definition}[Power Quantity]
  \label{def:physics:phyx-mini-0479:phyxminiproblems-problemphyxmini0479-powerquantity}
  \lean{PhyXMiniProblems.ProblemPhyXMini0479.PowerQuantity}
  Mechanical power carrying dimension M L² T⁻³
\end{definition}
\begin{proof}
  This is the declared model definition of Power Quantity.
\end{proof}
\begin{definition}[pressure In Pascals]
  \label{def:physics:phyx-mini-0479:phyxminiproblems-problemphyxmini0479-pressureinpascals}
  \lean{PhyXMiniProblems.ProblemPhyXMini0479.pressureInPascals}
  Read a physical pressure in pascals
\end{definition}
\begin{proof}
  This is the declared model definition of pressure In Pascals.
\end{proof}
\begin{definition}[pressure In Kilopascals]
  \label{def:physics:phyx-mini-0479:phyxminiproblems-problemphyxmini0479-pressureinkilopascals}
  \lean{PhyXMiniProblems.ProblemPhyXMini0479.pressureInKilopascals}
  \uses{def:physics:phyx-mini-0479:phyxminiproblems-problemphyxmini0479-pressureinpascals}
  Read a physical pressure in the kilopascals printed on the vertical axis
\end{definition}
\begin{proof}
  This is the declared model definition of pressure In Kilopascals.
\end{proof}
\begin{definition}[volume In Cubic Metres]
  \label{def:physics:phyx-mini-0479:phyxminiproblems-problemphyxmini0479-volumeincubicmetres}
  \lean{PhyXMiniProblems.ProblemPhyXMini0479.volumeInCubicMetres}
  \uses{def:physics:phyx-mini-0479:phyxminiproblems-problemphyxmini0479-volumequantity}
  Read a physical volume in coherent-SI cubic metres
\end{definition}
\begin{proof}
  This is the declared model definition of volume In Cubic Metres.
\end{proof}
\begin{definition}[volume In Cubic Centimetres]
  \label{def:physics:phyx-mini-0479:phyxminiproblems-problemphyxmini0479-volumeincubiccentimetres}
  \lean{PhyXMiniProblems.ProblemPhyXMini0479.volumeInCubicCentimetres}
  \uses{def:physics:phyx-mini-0479:phyxminiproblems-problemphyxmini0479-volumequantity, def:physics:phyx-mini-0479:phyxminiproblems-problemphyxmini0479-volumeincubicmetres}
  Read a physical volume in the cubic centimetres printed on the horizontal axis
\end{definition}
\begin{proof}
  This is the declared model definition of volume In Cubic Centimetres.
\end{proof}
\begin{definition}[energy In Joules]
  \label{def:physics:phyx-mini-0479:phyxminiproblems-problemphyxmini0479-energyinjoules}
  \lean{PhyXMiniProblems.ProblemPhyXMini0479.energyInJoules}
  Read signed heat, work, or internal-energy change in joules
\end{definition}
\begin{proof}
  This is the declared model definition of energy In Joules.
\end{proof}
\begin{definition}[temperature In Kelvins]
  \label{def:physics:phyx-mini-0479:phyxminiproblems-problemphyxmini0479-temperatureinkelvins}
  \lean{PhyXMiniProblems.ProblemPhyXMini0479.temperatureInKelvins}
  Read an absolute temperature in kelvins from its stated storage unit
\end{definition}
\begin{proof}
  This is the declared model definition of temperature In Kelvins.
\end{proof}
\begin{definition}[rotation Rate In Revolutions Per Second]
  \label{def:physics:phyx-mini-0479:phyxminiproblems-problemphyxmini0479-rotationrateinrevolutionspersecond}
  \lean{PhyXMiniProblems.ProblemPhyXMini0479.rotationRateInRevolutionsPerSecond}
  \uses{def:physics:phyx-mini-0479:phyxminiproblems-problemphyxmini0479-rotationratequantity}
  Read shaft speed in revolutions per SI second
\end{definition}
\begin{proof}
  This is the declared model definition of rotation Rate In Revolutions Per Second.
\end{proof}
\begin{definition}[rotation Rate In Revolutions Per Minute]
  \label{def:physics:phyx-mini-0479:phyxminiproblems-problemphyxmini0479-rotationrateinrevolutionsperminute}
  \lean{PhyXMiniProblems.ProblemPhyXMini0479.rotationRateInRevolutionsPerMinute}
  \uses{def:physics:phyx-mini-0479:phyxminiproblems-problemphyxmini0479-rotationratequantity, def:physics:phyx-mini-0479:phyxminiproblems-problemphyxmini0479-rotationrateinrevolutionspersecond}
  Read shaft speed in revolutions per minute
\end{definition}
\begin{proof}
  This is the declared model definition of rotation Rate In Revolutions Per Minute.
\end{proof}
\begin{definition}[cycle Rate In Hertz]
  \label{def:physics:phyx-mini-0479:phyxminiproblems-problemphyxmini0479-cyclerateinhertz}
  \lean{PhyXMiniProblems.ProblemPhyXMini0479.cycleRateInHertz}
  \uses{def:physics:phyx-mini-0479:phyxminiproblems-problemphyxmini0479-cycleratequantity}
  Read thermodynamic cycle rate in cycles per SI second
\end{definition}
\begin{proof}
  This is the declared model definition of cycle Rate In Hertz.
\end{proof}
\begin{definition}[power In Watts]
  \label{def:physics:phyx-mini-0479:phyxminiproblems-problemphyxmini0479-powerinwatts}
  \lean{PhyXMiniProblems.ProblemPhyXMini0479.powerInWatts}
  \uses{def:physics:phyx-mini-0479:phyxminiproblems-problemphyxmini0479-powerquantity}
  Read mechanical power in coherent-SI watts
\end{definition}
\begin{proof}
  This is the declared model definition of power In Watts.
\end{proof}
\begin{definition}[Gas Model]
  \label{def:physics:phyx-mini-0479:phyxminiproblems-problemphyxmini0479-gasmodel}
  \lean{PhyXMiniProblems.ProblemPhyXMini0479.GasModel}
  The working-substance model stated in the problem
\end{definition}
\begin{proof}
  This is the declared model definition of Gas Model.
\end{proof}
\begin{definition}[Thermodynamic Device Role]
  \label{def:physics:phyx-mini-0479:phyxminiproblems-problemphyxmini0479-thermodynamicdevicerole}
  \lean{PhyXMiniProblems.ProblemPhyXMini0479.ThermodynamicDeviceRole}
  The thermodynamic role played by the cyclic device
\end{definition}
\begin{proof}
  This is the declared model definition of Thermodynamic Device Role.
\end{proof}
\begin{definition}[State Label]
  \label{def:physics:phyx-mini-0479:phyxminiproblems-problemphyxmini0479-statelabel}
  \lean{PhyXMiniProblems.ProblemPhyXMini0479.StateLabel}
  The three numbered equilibrium states printed in the diagram
\end{definition}
\begin{proof}
  This is the declared model definition of State Label.
\end{proof}
\begin{definition}[Cycle Leg]
  \label{def:physics:phyx-mini-0479:phyxminiproblems-problemphyxmini0479-cycleleg}
  \lean{PhyXMiniProblems.ProblemPhyXMini0479.CycleLeg}
  The three directed arrows comprising one traversal of the cycle
\end{definition}
\begin{proof}
  This is the declared model definition of Cycle Leg.
\end{proof}
\begin{definition}[initial State]
  \label{def:physics:phyx-mini-0479:phyxminiproblems-problemphyxmini0479-cycleleg-initialstate}
  \lean{PhyXMiniProblems.ProblemPhyXMini0479.CycleLeg.initialState}
  \uses{def:physics:phyx-mini-0479:phyxminiproblems-problemphyxmini0479-statelabel, def:physics:phyx-mini-0479:phyxminiproblems-problemphyxmini0479-cycleleg}
  Initial state of a directed cycle leg
\end{definition}
\begin{proof}
  This is the declared model definition of initial State.
\end{proof}
\begin{definition}[final State]
  \label{def:physics:phyx-mini-0479:phyxminiproblems-problemphyxmini0479-cycleleg-finalstate}
  \lean{PhyXMiniProblems.ProblemPhyXMini0479.CycleLeg.finalState}
  \uses{def:physics:phyx-mini-0479:phyxminiproblems-problemphyxmini0479-statelabel, def:physics:phyx-mini-0479:phyxminiproblems-problemphyxmini0479-cycleleg}
  Final state of a directed cycle leg
\end{definition}
\begin{proof}
  This is the declared model definition of final State.
\end{proof}
\begin{definition}[Process Kind]
  \label{def:physics:phyx-mini-0479:phyxminiproblems-problemphyxmini0479-processkind}
  \lean{PhyXMiniProblems.ProblemPhyXMini0479.ProcessKind}
  Thermodynamic constraint holding along a process leg
\end{definition}
\begin{proof}
  This is the declared model definition of Process Kind.
\end{proof}
\begin{definition}[Path Geometry]
  \label{def:physics:phyx-mini-0479:phyxminiproblems-problemphyxmini0479-pathgeometry}
  \lean{PhyXMiniProblems.ProblemPhyXMini0479.PathGeometry}
  Geometric appearance of a process leg in the pressure--volume plane
\end{definition}
\begin{proof}
  This is the declared model definition of Path Geometry.
\end{proof}
\begin{definition}[Axis Quantity]
  \label{def:physics:phyx-mini-0479:phyxminiproblems-problemphyxmini0479-axisquantity}
  \lean{PhyXMiniProblems.ProblemPhyXMini0479.AxisQuantity}
  Physical quantity assigned to a plot axis
\end{definition}
\begin{proof}
  This is the declared model definition of Axis Quantity.
\end{proof}
\begin{definition}[Axis Unit]
  \label{def:physics:phyx-mini-0479:phyxminiproblems-problemphyxmini0479-axisunit}
  \lean{PhyXMiniProblems.ProblemPhyXMini0479.AxisUnit}
  Unit explicitly printed beside a plot axis
\end{definition}
\begin{proof}
  This is the declared model definition of Axis Unit.
\end{proof}
\begin{definition}[Figure Label]
  \label{def:physics:phyx-mini-0479:phyxminiproblems-problemphyxmini0479-figurelabel}
  \lean{PhyXMiniProblems.ProblemPhyXMini0479.FigureLabel}
  Text and tick labels visible in the primary bitmap
\end{definition}
\begin{proof}
  This is the declared model definition of Figure Label.
\end{proof}
\begin{definition}[Thermodynamic State]
  \label{def:physics:phyx-mini-0479:phyxminiproblems-problemphyxmini0479-thermodynamicstate}
  \lean{PhyXMiniProblems.ProblemPhyXMini0479.ThermodynamicState}
  \uses{def:physics:phyx-mini-0479:phyxminiproblems-problemphyxmini0479-volumequantity}
  A physical equilibrium state of the working gas
\end{definition}
\begin{proof}
  This is the declared model definition of Thermodynamic State.
\end{proof}
\begin{definition}[Pressure Volume Figure]
  \label{def:physics:phyx-mini-0479:phyxminiproblems-problemphyxmini0479-pressurevolumefigure}
  \lean{PhyXMiniProblems.ProblemPhyXMini0479.PressureVolumeFigure}
  \uses{def:physics:phyx-mini-0479:phyxminiproblems-problemphyxmini0479-statelabel, def:physics:phyx-mini-0479:phyxminiproblems-problemphyxmini0479-cycleleg, def:physics:phyx-mini-0479:phyxminiproblems-problemphyxmini0479-pathgeometry, def:physics:phyx-mini-0479:phyxminiproblems-problemphyxmini0479-axisquantity, def:physics:phyx-mini-0479:phyxminiproblems-problemphyxmini0479-axisunit, def:physics:phyx-mini-0479:phyxminiproblems-problemphyxmini0479-figurelabel}
  Qualitative and labelled information transcribed from image 479.png. This structure records no work, power, or answer choice
\end{definition}
\begin{proof}
  This is the declared model definition of Pressure Volume Figure.
\end{proof}
\begin{definition}[Ideal Gas Heat Engine Cycle]
  \label{def:physics:phyx-mini-0479:phyxminiproblems-problemphyxmini0479-idealgasheatenginecycle}
  \lean{PhyXMiniProblems.ProblemPhyXMini0479.IdealGasHeatEngineCycle}
  \uses{def:physics:phyx-mini-0479:phyxminiproblems-problemphyxmini0479-rotationratequantity, def:physics:phyx-mini-0479:phyxminiproblems-problemphyxmini0479-cycleratequantity, def:physics:phyx-mini-0479:phyxminiproblems-problemphyxmini0479-powerquantity, def:physics:phyx-mini-0479:phyxminiproblems-problemphyxmini0479-gasmodel, def:physics:phyx-mini-0479:phyxminiproblems-problemphyxmini0479-thermodynamicdevicerole, def:physics:phyx-mini-0479:phyxminiproblems-problemphyxmini0479-statelabel, def:physics:phyx-mini-0479:phyxminiproblems-problemphyxmini0479-cycleleg, def:physics:phyx-mini-0479:phyxminiproblems-problemphyxmini0479-processkind, def:physics:phyx-mini-0479:phyxminiproblems-problemphyxmini0479-thermodynamicstate, def:physics:phyx-mini-0479:phyxminiproblems-problemphyxmini0479-pressurevolumefigure}
  Independent physical observables for one closed engine cycle. The requested power is an independent dimensionful field; its relation to work and cycle rate is imposed only by the general governing law below
\end{definition}
\begin{proof}
  This is the declared model definition of Ideal Gas Heat Engine Cycle.
\end{proof}
\begin{definition}[net Work Done By Gas In Joules]
  \label{def:physics:phyx-mini-0479:phyxminiproblems-problemphyxmini0479-networkdonebygasinjoules}
  \lean{PhyXMiniProblems.ProblemPhyXMini0479.netWorkDoneByGasInJoules}
  \uses{def:physics:phyx-mini-0479:phyxminiproblems-problemphyxmini0479-energyinjoules, def:physics:phyx-mini-0479:phyxminiproblems-problemphyxmini0479-idealgasheatenginecycle}
  Net work done by the gas in one traversal of 1 → 2 → 3 → 1
\end{definition}
\begin{proof}
  This is the declared model definition of net Work Done By Gas In Joules.
\end{proof}
\begin{definition}[Matches Problem Statement]
  \label{def:physics:phyx-mini-0479:phyxminiproblems-problemphyxmini0479-matchesproblemstatement}
  \lean{PhyXMiniProblems.ProblemPhyXMini0479.MatchesProblemStatement}
  \uses{def:physics:phyx-mini-0479:phyxminiproblems-problemphyxmini0479-rotationrateinrevolutionsperminute, def:physics:phyx-mini-0479:phyxminiproblems-problemphyxmini0479-idealgasheatenginecycle}
  Data and idealized operating assumptions from the prose. The shaft makes one traversal of the displayed thermodynamic cycle per revolution; this explicit coupling is what permits the stated rpm to determine a cycle frequency. No numerical work or power occurs here
\end{definition}
\begin{proof}
  This is the declared model definition of Matches Problem Statement.
\end{proof}
\begin{definition}[Matches Primary Pressure Volume Figure]
  \label{def:physics:phyx-mini-0479:phyxminiproblems-problemphyxmini0479-matchesprimarypressurevolumefigure}
  \lean{PhyXMiniProblems.ProblemPhyXMini0479.MatchesPrimaryPressureVolumeFigure}
  \uses{def:physics:phyx-mini-0479:phyxminiproblems-problemphyxmini0479-pressureinkilopascals, def:physics:phyx-mini-0479:phyxminiproblems-problemphyxmini0479-volumeincubiccentimetres, def:physics:phyx-mini-0479:phyxminiproblems-problemphyxmini0479-temperatureinkelvins, def:physics:phyx-mini-0479:phyxminiproblems-problemphyxmini0479-statelabel, def:physics:phyx-mini-0479:phyxminiproblems-problemphyxmini0479-cycleleg, def:physics:phyx-mini-0479:phyxminiproblems-problemphyxmini0479-figurelabel, def:physics:phyx-mini-0479:phyxminiproblems-problemphyxmini0479-idealgasheatenginecycle}
  Exact transcription of the pressure--volume bitmap. The image gives the three directed legs, the process geometries, the state-1 and state-2 pressure and volume coordinates, the state-3 volume coordinate, and the 300 K isotherm through states 3 and 1. It does not give state 3 a numerical pressure and says nothing about net work or power
\end{definition}
\begin{proof}
  This is the declared model definition of Matches Primary Pressure Volume Figure.
\end{proof}
\begin{definition}[Has Physical Thermodynamic Parameters]
  \label{def:physics:phyx-mini-0479:phyxminiproblems-problemphyxmini0479-hasphysicalthermodynamicparameters}
  \lean{PhyXMiniProblems.ProblemPhyXMini0479.HasPhysicalThermodynamicParameters}
  \uses{def:physics:phyx-mini-0479:phyxminiproblems-problemphyxmini0479-pressureinpascals, def:physics:phyx-mini-0479:phyxminiproblems-problemphyxmini0479-volumeincubicmetres, def:physics:phyx-mini-0479:phyxminiproblems-problemphyxmini0479-temperatureinkelvins, def:physics:phyx-mini-0479:phyxminiproblems-problemphyxmini0479-rotationrateinrevolutionspersecond, def:physics:phyx-mini-0479:phyxminiproblems-problemphyxmini0479-cyclerateinhertz, def:physics:phyx-mini-0479:phyxminiproblems-problemphyxmini0479-powerinwatts, def:physics:phyx-mini-0479:phyxminiproblems-problemphyxmini0479-idealgasheatenginecycle}
  Positivity and nondegeneracy conditions for the physical branch
\end{definition}
\begin{proof}
  This is the declared model definition of Has Physical Thermodynamic Parameters.
\end{proof}
\begin{definition}[Satisfies Monatomic Ideal Gas Engine Laws]
  \label{def:physics:phyx-mini-0479:phyxminiproblems-problemphyxmini0479-satisfiesmonatomicidealgasenginelaws}
  \lean{PhyXMiniProblems.ProblemPhyXMini0479.SatisfiesMonatomicIdealGasEngineLaws}
  \uses{def:physics:phyx-mini-0479:phyxminiproblems-problemphyxmini0479-pressureinpascals, def:physics:phyx-mini-0479:phyxminiproblems-problemphyxmini0479-volumeincubicmetres, def:physics:phyx-mini-0479:phyxminiproblems-problemphyxmini0479-energyinjoules, def:physics:phyx-mini-0479:phyxminiproblems-problemphyxmini0479-temperatureinkelvins, def:physics:phyx-mini-0479:phyxminiproblems-problemphyxmini0479-rotationrateinrevolutionspersecond, def:physics:phyx-mini-0479:phyxminiproblems-problemphyxmini0479-cyclerateinhertz, def:physics:phyx-mini-0479:phyxminiproblems-problemphyxmini0479-powerinwatts, def:physics:phyx-mini-0479:phyxminiproblems-problemphyxmini0479-statelabel, def:physics:phyx-mini-0479:phyxminiproblems-problemphyxmini0479-cycleleg, def:physics:phyx-mini-0479:phyxminiproblems-problemphyxmini0479-idealgasheatenginecycle, def:physics:phyx-mini-0479:phyxminiproblems-problemphyxmini0479-networkdonebygasinjoules}
  Macroscopic governing laws for a fixed sample of monatomic ideal gas: the cross-multiplied ideal-gas law expresses pV/T = constant; monatomic internal energy is (3/2) pV, up to a state-independent datum; the first law is Q\_into = ΔU + W\_by; isobaric, isochoric, and reversible-isothermal boundary-work formulas hold; cycle rate is shaft rate times cycles per revolution; average output power is net work per cycle times cycles per second. These are general laws.
\end{definition}
\begin{proof}
  This is the declared model definition of Satisfies Monatomic Ideal Gas Engine Laws.
\end{proof}
\begin{lemma}[derived Cycle Readouts]
  \label{lem:physics:phyx-mini-0479:phyxminiproblems-problemphyxmini0479-derivedcyclereadouts}
  \lean{PhyXMiniProblems.ProblemPhyXMini0479.derivedCycleReadouts}
  \uses{def:physics:phyx-mini-0479:phyxminiproblems-problemphyxmini0479-pressureinkilopascals, def:physics:phyx-mini-0479:phyxminiproblems-problemphyxmini0479-energyinjoules, def:physics:phyx-mini-0479:phyxminiproblems-problemphyxmini0479-cyclerateinhertz, def:physics:phyx-mini-0479:phyxminiproblems-problemphyxmini0479-idealgasheatenginecycle, def:physics:phyx-mini-0479:phyxminiproblems-problemphyxmini0479-networkdonebygasinjoules, def:physics:phyx-mini-0479:phyxminiproblems-problemphyxmini0479-matchesproblemstatement, def:physics:phyx-mini-0479:phyxminiproblems-problemphyxmini0479-matchesprimarypressurevolumefigure, def:physics:phyx-mini-0479:phyxminiproblems-problemphyxmini0479-hasphysicalthermodynamicparameters, def:physics:phyx-mini-0479:phyxminiproblems-problemphyxmini0479-satisfiesmonatomicidealgasenginelaws}
  The ideal-gas and process laws determine the previously unread pressure and the cycle's work and frequency. In particular, the three leg works are 80 J, 0 J, and -40 log 3 J, giving 40 (2 - log 3) J per cycle at 10 cycles/s
\end{lemma}
\begin{proof}
  The conclusion follows from the governing relations and figure readouts named in the statement of derived Cycle Readouts.
\end{proof}
\begin{definition}[Answer Choice]
  \label{def:physics:phyx-mini-0479:phyxminiproblems-problemphyxmini0479-answerchoice}
  \lean{PhyXMiniProblems.ProblemPhyXMini0479.AnswerChoice}
  Labels of the four multiple-choice answers
\end{definition}
\begin{proof}
  This is the declared model definition of Answer Choice.
\end{proof}
\begin{definition}[displayed Power In Watts]
  \label{def:physics:phyx-mini-0479:phyxminiproblems-problemphyxmini0479-displayedpowerinwatts}
  \lean{PhyXMiniProblems.ProblemPhyXMini0479.displayedPowerInWatts}
  \uses{def:physics:phyx-mini-0479:phyxminiproblems-problemphyxmini0479-answerchoice}
  Power in watts printed beside each answer choice
\end{definition}
\begin{proof}
  This is the declared model definition of displayed Power In Watts.
\end{proof}
\begin{definition}[recorded Answer Choice]
  \label{def:physics:phyx-mini-0479:phyxminiproblems-problemphyxmini0479-recordedanswerchoice}
  \lean{PhyXMiniProblems.ProblemPhyXMini0479.recordedAnswerChoice}
  \uses{def:physics:phyx-mini-0479:phyxminiproblems-problemphyxmini0479-answerchoice}
  Answer label recorded by the source dataset
\end{definition}
\begin{proof}
  This is the declared model definition of recorded Answer Choice.
\end{proof}
\begin{definition}[Rounds To Displayed Power]
  \label{def:physics:phyx-mini-0479:phyxminiproblems-problemphyxmini0479-roundstodisplayedpower}
  \lean{PhyXMiniProblems.ProblemPhyXMini0479.RoundsToDisplayedPower}
  \uses{def:physics:phyx-mini-0479:phyxminiproblems-problemphyxmini0479-powerinwatts, def:physics:phyx-mini-0479:phyxminiproblems-problemphyxmini0479-idealgasheatenginecycle, def:physics:phyx-mini-0479:phyxminiproblems-problemphyxmini0479-answerchoice, def:physics:phyx-mini-0479:phyxminiproblems-problemphyxmini0479-displayedpowerinwatts}
  The physical output lies in the nearest-20-watt bin of a displayed choice
\end{definition}
\begin{proof}
  This is the declared model definition of Rounds To Displayed Power.
\end{proof}
\begin{definition}[Is Unique Closest Displayed Power]
  \label{def:physics:phyx-mini-0479:phyxminiproblems-problemphyxmini0479-isuniqueclosestdisplayedpower}
  \lean{PhyXMiniProblems.ProblemPhyXMini0479.IsUniqueClosestDisplayedPower}
  \uses{def:physics:phyx-mini-0479:phyxminiproblems-problemphyxmini0479-powerinwatts, def:physics:phyx-mini-0479:phyxminiproblems-problemphyxmini0479-idealgasheatenginecycle, def:physics:phyx-mini-0479:phyxminiproblems-problemphyxmini0479-answerchoice, def:physics:phyx-mini-0479:phyxminiproblems-problemphyxmini0479-displayedpowerinwatts}
  One displayed power is strictly closer than every competing choice
\end{definition}
\begin{proof}
  This is the declared model definition of Is Unique Closest Displayed Power.
\end{proof}
% --- Archon declaration topology end ---
% --- Archon physics formalization source end ---
